\chapter{The Sieve Census}\label{appB:census}

This appendix records, for reference, the combinatorial data underlying the
executable topology closures discussed in Chapter~11. Every count in this
appendix is reproduced from the modality-composition verifier's registered
output or is directly derivable from the definitions in Chapters~7 and~9; none
is an independent claim. The governing status remains \textsc{draft / gate 1
research candidate / ring 4 / not licensed}, and exhaustive enumeration over one
sixteen-object site is corroboration of the finite proofs, not proof-assistant
verification or independent review.

\section{The underlying poset}\label{appB:poset}

Recall from Chapter~7 that $\Omega=\{(0,0),(0,1),(1,0),(1,1)\}$ and
$L=\mathcal{P}(\Omega)$ ordered by inclusion. We abbreviate an atom $(i,j)$ by
the two-digit string $ij$, so the four atoms are $00,01,10,11$ and a context is
written as the concatenation of its atoms in lexicographic order (the empty
context is written $\varnothing$).

\begin{longtable}{@{}lll@{}}
\caption{The sixteen contexts of $L$ by cardinality.}\label{appB:contexts}\\
\toprule
Size & Count & Contexts \\
\midrule
\endfirsthead
\toprule
Size & Count & Contexts \\
\midrule
\endhead
\bottomrule
\endfoot
$0$ & $1$ & $\varnothing$ \\
$1$ & $4$ & $00$, $01$, $10$, $11$ \\
$2$ & $6$ & $00\,01$ ($=A_0$), $10\,11$ ($=A_1$), $00\,10$ ($=R_0$), $01\,11$ ($=R_1$), $00\,11$, $01\,10$ \\
$3$ & $4$ & $00\,01\,10$, $00\,01\,11$, $00\,10\,11$, $01\,10\,11$ \\
$4$ & $1$ & $00\,01\,10\,11$ ($=\Omega$) \\
\bottomrule
\end{longtable}

The two size-two contexts $00\,11$ and $01\,10$ are the \emph{diagonal}
doubletons: each contains exactly one atom from each ACBE block and one from
each RPM block, so it contains no complete block of either partition. The four
block doubletons $A_0,A_1,R_0,R_1$ each contain exactly one complete block.

\section{Arrow count}\label{appB:arrows}

\begin{proposition}\label{appB:arrowcount}
The inclusion poset $L$ has exactly $81$ arrows (ordered pairs $U\subseteq V$).
\end{proposition}

\begin{proof}
An ordered pair $U\subseteq V$ is determined by independently assigning each of
the four atoms to one of three roles: in neither $U$ nor $V$; in $V$ but not
$U$; or in both. These roles are mutually exclusive and exhaustive because
$U\subseteq V$ forbids the fourth combination (in $U$ but not $V$). Hence the
count is $3^4=81$, matching the verifier's reported arrow total.
\end{proof}

\section{Sieve and covering-sieve totals}\label{appB:sieves}

The modality-composition verifier enumerates every sieve on every context and
saturates the Grothendieck axioms to a fixed point. Its registered totals are
reproduced in Table~\ref{appB:sievetotals}. The two presentations --- the join
$J_A\vee J_R$ and the singleton topology $J_{\mathrm{singleton}}$ of
Chapter~9 --- yield identical covering-sieve counts, corroborating the join
identification (NT2-4).

\begin{longtable}{@{}lr@{}}
\caption{Registered totals for the sixteen-context site: cardinalities,
sieve and covering-sieve counts, and the composition verifier's fixture
and restriction-square counts.}\label{appB:sievetotals}\\
\toprule
Quantity & Count \\
\midrule
\endfirsthead
\toprule
Quantity & Count \\
\midrule
\endhead
\bottomrule
\endfoot
Atoms & $4$ \\
Contexts & $16$ \\
Inclusion arrows & $81$ \\
Sieves (all contexts) & $298$ \\
Covering sieves, $J_A\vee J_R$ presentation & $168$ \\
Covering sieves, $J_{\mathrm{singleton}}$ presentation & $168$ \\
Executable fixtures (composition verifier) & $18$ \\
Restriction squares checked & $3012$ \\
\bottomrule
\end{longtable}

\begin{remark}\label{appB:sieveremark}
The count of $298$ sieves is the total across all sixteen contexts, not the
number on a single context: a context of size $k$ carries a lattice of sieves
whose size grows with $k$, and the census sums these. The equality of the two
$168$ covering-sieve counts is the executable face of the join identification
(the join lemma of Chapter~10): the topology generated jointly by the ACBE and
RPM families and the topology generated by the four singletons are literally the
same set of covering sieves, not merely isomorphic ones.
\end{remark}

\section{Block-completeness table}\label{appB:blocks}

Table~\ref{appB:blocktable} records, for each context, which complete ACBE and
RPM blocks it contains. This is the data that drives the witness presheaves of
Chapter~9: $F_A(V)$ has $2^{\alpha(V)}$ sections, where $\alpha(V)$ is the
number of complete ACBE blocks in $V$ (Definition of complete blocks,
Chapter~7), and symmetrically $|F_R(V)| = 2^{\rho(V)}$.

\begin{longtable}{@{}lll@{}}
\caption{Complete blocks contained in each context.}\label{appB:blocktable}\\
\toprule
Context & Complete ACBE blocks & Complete RPM blocks \\
\midrule
\endfirsthead
\toprule
Context & Complete ACBE blocks & Complete RPM blocks \\
\midrule
\endhead
\bottomrule
\endfoot
$\varnothing$ & none & none \\
$00$, $01$, $10$, $11$ & none & none \\
$00\,01$ ($A_0$) & $A_0$ & none \\
$10\,11$ ($A_1$) & $A_1$ & none \\
$00\,10$ ($R_0$) & none & $R_0$ \\
$01\,11$ ($R_1$) & none & $R_1$ \\
$00\,11$ & none & none \\
$01\,10$ & none & none \\
$00\,01\,10$ & $A_0$ & $R_0$ \\
$00\,01\,11$ & $A_0$ & $R_1$ \\
$00\,10\,11$ & $A_1$ & $R_0$ \\
$01\,10\,11$ & $A_1$ & $R_1$ \\
$\Omega$ & $A_0$, $A_1$ & $R_0$, $R_1$ \\
\bottomrule
\end{longtable}

\begin{remark}\label{appB:whyfour}
Only $\Omega$ contains both complete ACBE blocks, so only there does $F_A$ take
its full four sections. By transversality no RPM piece $V\cap R_j$ of any
context contains a complete ACBE block, so the RPM matching-family set for
$F_A$ is a one-element set at \emph{every} context (Appendix~C,
Finding~\ref{appC:fnd:novaluecrossing}); the gluing map therefore fails
injectivity wherever $|F_A(V)|\ge 2$, that is at the two ACBE blocks, the four
triples, and $\Omega$. The distinctness computation of Chapter~9 is carried out
at $\Omega$ because that is the context with the widest gap --- four sections
against one matching family --- not because it is the only failure.
\end{remark}
